\documentclass[11pt,a4paper]{article}

\usepackage[margin=2.5cm]{geometry}
\usepackage[T1]{fontenc}
\usepackage[utf8]{inputenc}
\usepackage{lmodern}
\usepackage{microtype}
\usepackage{booktabs}
\usepackage[hidelinks]{hyperref}

\title{Urban Outfinders\\[0.3em]\large Part 3: Ranking \& Filtering}
\author{Marc Mallol\and Èric Dürr\and Oriol Pont}
\date{Information Retrieval and Web Analytics, 2026}

\begin{document}
\maketitle

\section{Retrieval setup}
% Introduce the catalogue, searchable fields, preprocessing and the five Part 2
% queries. Explain AND matching, shared candidates and ties broken by product ID.

\section{Lexical and custom ranking}
\subsection{TF-IDF and cosine similarity}
% Summarize the Part 2 baseline and record its settings.

\subsection{BM25}
% Give the formula, IDF convention, k1 and b, and document length definition.

\subsection{Custom score}
% Justify the formula and feature choices, scaling, missing values and trade-offs.

\section{Word2Vec and cosine similarity}
\subsection{Text vectors}
% Describe the model, tokenization, averaging, out-of-vocabulary terms and
% zero-vector handling.

\subsection{Top 20 results per query}
% Insert a ranked table for each Part 2 test query. Include product IDs, titles
% and scores. Report the candidate count when fewer than 20 products match.

\section{Ranking comparison}
% Explain TF-IDF and BM25 rank changes and the pros and cons of both.
% Discuss the effects and trade-offs of the custom score and Word2Vec.
% If reporting metrics from Part 2 labels, explain the judgement pool and
% report coverage across rankers.

\section{Alternative text representations}
% Propose a representation and justify its pros and cons relative to averaged
% Word2Vec. Consider Doc2Vec and sentence embeddings.

\section{Conclusion}
% Summarize the observed ranking differences and which methods suited the queries.

\end{document}
